
\addchaptertocentry{Хураангуй}

\begin{center}
{\scshape\Large \univname\par}
{\scshape\large \facname\par}\vspace{0.5cm}
{\huge\textbf{{Хураангуй}} \par}
\bigskip
{\Large{\ttitle} \par}
\bigskip

{\normalsize \shortname \par}
\addressname
\end{center}

\textit{\textbf{Түлхүүр үгс: хиймэл оюун, цахим сургалт, дасан зохицох сургалт, ахиц хянах, чатбот, санал болгох систем}}
\bigskip

Энэхүү дипломын ажлын хүрээнд програмчлалын сургалтад зориулсан хиймэл оюунд суурилсан, шаталсан сургалтын веб системийг судалж, шинжилж, хөгжүүлэв. Ажлын зорилго нь хэрэглэгчийн мэдлэгийн түвшинг тодорхойлон түүнд тохирсон курс болон хичээл санал болгох, суралцах явцыг хянах, хичээлийн түвшинд болон системийн түвшинд хиймэл оюуны тусламж үзүүлэх цогц платформ бий болгоход оршино.

Системийн хүрээнд хэрэглэгчийн бүртгэл, нэвтрэлт, профайлын удирдлага, түвшин тогтоох шалгалт, түвшин ахих шалгалт, суралцах чиглэл, курс болон хичээлийн удирдлага, таалагдсан хичээлүүд, ахицын хураангуй, сэтгэгдлээ хуваалцах модуль, тоглоомын хэсэг зэрэг үндсэн боломжуудыг хэрэгжүүлсэн. Мөн хичээлийн агуулгад тулгуурлан хиймэл оюун ашиглаж хураангуй үүсгэх, асуулт бэлдэх, дадлагын даалгавар үүсгэх, чатботоор асуултад хариулах, сэтгэгдлийн хяналт хийх ухаалаг механизмуудыг нэгтгэсэн.

Хөгжүүлэлтийн хувьд урд талын хэсгийг React, TypeScript, Vite; арын хэсгийг Django болон Django REST Framework; өгөгдлийн сангийн түвшинд PostgreSQL; орчны удирдлагад Docker Compose ашиглав. JWT бүртгэл баталгаажуулалт, ахицын нэгтгэл, зөвлөмжийн логик, хиймэл оюуны нөөц шийдэл, нөөцийн хязгаарлалт бүхий код ажиллуулагч, аюулгүй байдлын орчны тохиргоо зэрэг инженерчлэлийн шийдлүүдийг хэрэгжүүлсэн.

Судалгаа, шинжилгээ, хөгжүүлэлт, туршилтын үр дүнд дасан зохицох сургалт, хиймэл оюуны тусламжтай дэмжлэг, ахицын шинжилгээ, сэтгэгдэл үлдээх харилцаа, тоглоомжуулалтыг нэгтгэсэн, цааш өргөтгөх боломжтой сургалтын системийн загвар боловсруулагдсан. Энэхүү ажил нь програмчлалын сургалтад хиймэл оюуныг практик түвшинд нэгтгэх, хувь хүнд тохирсон сургалтын урсгал бий болгох, орчин үеийн цогц веб технологиор шийдэл хэрэгжүүлэх жишиг суурь болсон.
